\begin{answer}
I agreed with Z on \textbf{1 of 5} questions (Larry Page) and disagreed on
\textbf{4 of 5}. E's response was longer than F's in all five pairs, and Z
picked E each time. I picked the shorter answer in four cases. I prefer
concise answers for simple questions and more detailed answers for complex
ones, especially when the extra detail helps the user take the next step.
Z's reasons were often relevant, but in this sample its choices did not
consistently match that preference.

\smallskip
\textbf{Disagreement example:} I picked F because it clearly asks for the
clue, the number of letters, and known letters. That is enough to start.
Z picked E, citing its examples, grid details, and explanation of the help
it could give. The original English responses are below. In this pair, the
blind labels were A for F and B for E.
\end{answer}

\medskip
\noindent\textbf{AlpacaEval question}\par
\begin{lstlisting}[language={},basicstyle=\small\ttfamily,breaklines=true,frame=none]
Hi, I'm trying to solve a crossword puzzle, but I've never done one of these before. Can you help me out?
\end{lstlisting}

\medskip
\noindent\textbf{Response E (blind label B)}\par
\begin{lstlisting}[language={},basicstyle=\small\ttfamily,breaklines=true,frame=none]
Absolutely -- I'd be happy to help.

If you want, send me any of the following:
- the crossword clue
- the number of letters
- any letters you already have
- the theme, if there is one
- a photo or typed layout of the grid

For example, you can say:
- **"Clue: 'Large deer,' 4 letters"**
- **"I have _ E _ _ for 'Large deer'"**

And I can help by:
- suggesting likely answers
- explaining *why* an answer fits
- helping with tricky clue types like puns, abbreviations, or wordplay
- narrowing options based on crossing letters

If you want, we can also do it step by step so it feels less overwhelming. Send me your first clue.
\end{lstlisting}

\medskip
\noindent\textbf{Response F (blind label A)}\par
\begin{lstlisting}[language={},basicstyle=\small\ttfamily,breaklines=true,frame=none]
Absolutely -- I can help.

If you send me:
- the clue
- the number of letters
- any letters you already have
- and whether it's **across** or **down** (if relevant)

...I can help you figure it out.

If you want, I can also explain **how crossword clues usually work** so you can get better at solving them yourself.
\end{lstlisting}

\medskip
\noindent\textbf{Judge Z}\par
\begin{lstlisting}[language={},basicstyle=\small\ttfamily,breaklines=true,frame=none]
B is more useful because it asks for the key details needed to solve a crossword, gives concrete examples of how to provide them, and explains the kinds of help it can offer, including themes, grid layout, and crossing-letter logic. A is good, but B is more actionable for a beginner.

WINNER: B
\end{lstlisting}
